\section{Discussion}
\label{sec:discussion}

\textbf{What we learned.} The largest effect in this report is not a difference
between models. It is the drop from training to unseen demonstrations, which
every model shows and which is not caused by drift over the session. \rev{Nor is
it caused by the test demonstrations looking unfamiliar, and it is built late in
training: for all but FastWAM, test error levels off early while training error
keeps falling.}
Fifty-eight demonstrations of one task from one session are not enough for any
of these models to generalise to a fifty-ninth.

Cropping the tactile images is not the lever. It changes almost nothing on any
model, and \rev{no model} over-weights the image edges that motivated it.
\new{Pretraining did not help as we used it: pi0.5, finetuned through a narrow
LoRA, plans worse than a small flow-matching policy trained from scratch, and
three times more training did not close the gap.}

Touch is used, but briefly, \new{and removing the fingertips altogether made
ACT slightly better on held-out demonstrations. On a task this consistent, the
overhead camera and the joint positions may already say what the next commands
should be.} The two world action models plan the near future best, though
DreamZero barely reads its cameras to do so.

\textbf{Next steps.} \textbf{Collect more varied demonstrations} (sessions,
garments, starting positions), the most direct attack on the generalisation
gap. \textbf{Test on the robot}, since agreement with recorded commands is not
folding the garment. \new{\textbf{Finetune pi0.5 fully, or with LoRA on every
layer,} as the published recipes do.} \textbf{Score motion, not whole frames}, which frame-averaged PSNR ignores.
